% A local obstruction for the proposed Refl deletion.
\section{Why a neutral equality does not replace reflexivity}
\label{sec:reflexivity-obstruction}

A tempting further reduction is to remove Refl and retain only the
neutral equality $H=(0=0)$ used by ordinary assumption clauses.
The following obstruction shows that this loses even positive-numeral
reflexivity.

Let $\mathsf{CE5}^{0}$ have the same raw terms, formulas, contexts,
$\Phi$, and rules U, E, I, and Eval as $\mathsf{CE5}$, but replace
the Refl schema by the single axiom $0=0$. Its axiom is available in
every context.

\begin{theorem}[The neutral equality is insufficient]
\label{thm:reflexivity-obstruction}
The weakened calculus $\mathsf{CE5}^{0}$ does not derive $S0=S0$
from the empty context. In particular, it is not consequence-equivalent
to $\mathsf{CE5}$ on their common raw language.
\end{theorem}
\begin{proof}
Evaluate all numerical terms by their usual operations on the standard
natural numbers, including ordinary initialized iteration. However,
assign to the \emph{syntactic equation atom} $s=t$ the truth condition
\[
 \llbracket s=t\rrbracket_\nu
 \quad\Longleftrightarrow\quad
 \bigl(\llbracket s\rrbracket_\nu=0
       \ \hbox{and}\ \llbracket t\rrbracket_\nu=0\bigr).
                                                        \tag{Z}
\]
Thus both sides must denote zero. Interpret $D$ and clauses by
\[
 \llbracket D_{x,y}B\rrbracket
  =\exists x\bigl(\llbracket B(x,0)\rrbracket
                  \to\forall y\,\llbracket B(x,y)\rrbracket\bigr),
\]
\[
 \llbracket[A;y.B]\rrbracket
  =\llbracket A\rrbracket\to\forall y\,\llbracket B(y)\rrbracket,
\]
using ordinary two-valued truth. An empty context is true.

We verify preservation for every primitive rule. The axiom $0=0$
satisfies (Z). For Eval, the truth of its two premises implies
\[
 \llbracket s\rrbracket
 =\llbracket t\rrbracket
 =\llbracket b\rrbracket
 =\llbracket r\rrbracket=0.
\]
In the root row of $\Phi$ its value is $t$, so the conclusion $r=t$
satisfies (Z). In the zero-iterator row its value is again $t$, with
the same conclusion. In the successor row write
$b=I(Sn,a;u.h)$ and $s=I(n,a;u.h)$. The ordinary numerical evaluation
satisfies
\[
 \begin{aligned}
 \llbracket h[t/u]\rrbracket
 &=\llbracket h[s/u]\rrbracket\\
 &=\llbracket I(Sn,a;u.h)\rrbracket
  =\llbracket b\rrbracket=0.
 \end{aligned}
\]
The first equality follows because $s$ and $t$ both have ordinary
numerical value zero; the second is ordinary iterator computation.
Together with $\llbracket r\rrbracket=0$ this makes the Eval
conclusion true.

U is valid implication/universal specialization. E is valid existential
elimination: its fresh parameter may be assigned the existential
witness, and the conclusion and retained context do not depend on
that parameter.

For I in case mode, fix an assignment satisfying the retained context
and assume the matrix's zero instance at the value of its witness.
Validity of the premise gives every successor instance. Every standard
natural number is zero or a successor, so all instances hold, proving
the desired conditional closure and its existential conclusion.

For I in progress mode, its witness is literally zero. The valid
premise says, for every standard $n$, that the truth of $B(0,n)$
implies the truth of $B(0,Sn)$, under the fixed retained context.
If $B(0,0)$ holds, ordinary metatheoretic induction on standard natural
numbers gives every instance. If the base does not hold, the
conditional closure is already true. This argument is valid for
the interpretation (Z); it does not require equation atoms to mean
ordinary equality. All fresh-variable conditions ensure that the
retained context and witness stay fixed as $n$ varies.

Induction on finite derivations therefore shows that every closed
theorem of $\mathsf{CE5}^{0}$ is true under this interpretation.
But $S0$ has ordinary value $1$, so (Z) makes $S0=S0$ false.
\end{proof}

The interpretation in this proof is a counterinterpretation of a
syntactic inference system. It is deliberately \emph{not} a
first-order arithmetic structure whose equality symbol is interpreted
as identity, nor is it the proposed arithmetic interpretation of the
new theory. Such a counterinterpretation suffices to establish
nonderivability: each permitted inference preserves its truth
condition, while the desired conclusion fails it.

If Refl is deleted without even adding $0=0$, the still simpler
valuation making every equation atom false gives the analogous
obstruction.

\begin{corollary}[No finite collection of closed equation seeds]
Delete Refl and add any finite collection of closed equation axioms
that are true under ordinary standard-natural-number evaluation.
Keeping Eval, U, E, and I unchanged still does not recover universal
reflexivity.
\end{corollary}
\begin{proof}
Let $A\subseteq\mathbb N$ be the finite set of numerical values of
the supplied axioms' sides. Both sides of each axiom have the same
value, since the axioms were assumed ordinarily true. Replace (Z) by
\[
 \llbracket s=t\rrbracket_\nu
 \quad\Longleftrightarrow\quad
 \llbracket s\rrbracket_\nu
   =\llbracket t\rrbracket_\nu\in A.
\]
Every seed is true. In each Eval row, actual equality of its first
premise's values and ordinary iterator computation imply that
$\Phi(s,t,b)$ has the ordinary value of $b$. Its second premise
then says that this value equals the value of $r$ and belongs to $A$,
so the conclusion is true. The previous U/E argument and the case
and progress arguments for I apply to this arbitrary atomic predicate
on standard $\mathbb N$ without change.

Choose $k\notin A$. The closed numeral equation
$S^k0=S^k0$ is false in this rule-preserving interpretation and hence
is not derivable. This is a local obstruction for the fixed remaining
rules and closed equation seeds; it makes no claim about other axiom
forms or different calculi.
\end{proof}

\paragraph{A retained reflexivity clause is an axiom.}
The permanent clause
\[
 [\,0=0;\ y.\,y=y\,]
\]
together with a proof of $0=0$ does recover Refl: U at target $t$
gives $t=t$. That clause has the ordinary reading
$(0=0)\to\forall y(y=y)$ and fails the counterinterpretation above.
Initializing a checker with this clause and a neutral equality proof
can remove an instruction tag, but it has relocated the missing
universal reflexivity axiom into the initialization. Both the axiom
and its validation cost must remain in the complete specification.
It is not an additional mathematical rule deletion, and it does
not justify a four-pattern equivalence claim for the four remaining
rules alone.

